\documentclass[12pt]{hw}

% ====================================================================
% IMPORTANT INFORMATION TO FILL OUT:
% ====================================================================

% ENTER YOUR NAME HERE:
\renewcommand{\student}{ YOUR NAME HERE (uniqname here) }

% LIST ALL COLLABORATORS HERE:
\renewcommand{\collaborators}{Name/Exercise Number(s), Name/Exercise
  Numbers(s)}

% ====================================================================
% ====================================================================

\setcounter{sheet}{1} % Homework number here

\begin{document}

\maketitle

\begin{center}
  Due: September 12, 2024 (at 11:59 PM US Eastern Time) \\[10pt]
  \hrule height 1pt
\end{center}

% ====================================================================
% ====================================================================

\begin{exercise}
  \textbf{(10 points)} Recall that an encryption scheme is formally
  defined by algorithms $\skcgen, \skcenc$, and $\skcdec$
  (respectively, the key generator, encryption, and decryption),
  together with a message space, cipher space, and key space.  Give
  formal specifications for the shift cipher and the
  Vigen\`ere cipher. For the latter, you may assume that the key has a length 5.

  Vigen\`ere cipher is covered in the second discussion lecture. You can also find it in Section 1.1.3 from Vipul Goyal's lecture notes (see Canvas for the link). 
\end{exercise}

\begin{answer}
  Solution here.
\end{answer}

% ====================================================================
% ====================================================================

\begin{exercise}
  \textbf{(10 points)} 
You are given three identical cards, with the exception that both sides of the first card are
red, both sides of the second card are black, and one side of the third card is red and the
other black. The cards are shuffled, and one card is randomly drawn and put down on
the ground. If the upper side of the chosen card is red, what is the probability that the
card on the ground has distinct colors on its sides?
\end{exercise}

\begin{answer}
  Solution here.
\end{answer}

% ====================================================================
% ====================================================================

\begin{exercise}
  \textbf{(20 points)} 
\begin{enumerate}[(a)]
\item 
Suppose that for some private-key encryption scheme,
the ciphertext always starts with the letter ``c'', 
can this be a perfectly secure encryption?
Please explain why.
\item 
Suppose that for some private-key encryption scheme,
the key always starts with the letter ``c''; 
can this be a perfectly secure encryption?
Please explain why.
\item 
Suppose a message is a string containing letters from 
``a'' to ``z''. If the ciphertext always contains the first letter of the message, can it be perfectly secure encryption? 
Please explain why.
\end{enumerate}
\end{exercise}

\begin{answer}
  Solution here.
  \begin{enumerate}[(a)]
  \item

  \item

  \item
  \end{enumerate}
\end{answer}

% ====================================================================
% ====================================================================

\begin{exercise}
  \textbf{(20 points)} Let $\mathbb{Z}_n$ denote integers modulo $n$,
and let $\mathbb{Z}_n \setminus \{0\}$ denote all the elements of $\mathbb{Z}_n$ except $0$ (e.g. $\mathbb{Z}_3 = \{0, 1, 2\}$ and $\mathbb{Z}_3 \setminus \{0\} = \{1, 2\}$). 
% Definition
For some fixed message length $\ell$,
let keyspace $\mathcal{K}$, message space $\mathcal{M}$, and cipertext space $\mathcal{C}$ all be $(\mathbb{Z}_{n} \setminus \{0\})^\ell$ (i.e. vectors of length $\ell$ over elements from $\mathbb{Z}_{n} \setminus \{0\}$).
We partially define the following encryption scheme over $\mathbb{Z}_n \setminus \{0\}$:
\begin{align*}
    \mathsf{Gen}(1^\ell) &: (k_1, \dots, k_\ell) \sample (\mathbb{Z}_{n} \setminus \{0\})^\ell\\
    \mathsf{Enc}(k, m) &: (k_1 \cdot m_1, k_2 \cdot m_2, \dots, k_\ell \cdot m_\ell)\mod n
\end{align*}
\begin{enumerate}[(a)]
    \item Define a decryption algorithm for the above encryption scheme over $\mathbb{Z}_3 \setminus \{0\}$
    (i.e. $\mathcal{K} = \mathcal{M} = \mathcal{C} = \mathbb{Z}_{3} \setminus \{0\}$), 
    and prove its correctness.
    \item 
    Show that the encryption scheme over $\mathbb{Z}_{3} \setminus \{0\}$ is perfectly secret.
    \item Suppose we define the encryption scheme over $\mathbb{Z}_4 \setminus \{0\}=\{1,2,3\}$. 
    Is this a valid encryption scheme?
\end{enumerate}
\end{exercise}

\begin{answer}
  Solution here.

  \begin{enumerate}[(a)]
  \item

  \item

  \item
  \end{enumerate}

\end{answer}

% ====================================================================
% ====================================================================

\begin{exercise}
  \textbf{(20 points)} Prove or give a counter example: An encryption scheme with a
  finite message space $\mathcal{M}$ is perfectly secret if and only
  if for every probability distribution $\mathcal{D}$ over
  $\mathcal{M}$ and every $c_0, c_1 \in \mathcal{C}$ over the cipher
  space $\mathcal{C}$, we have $\Pr[C = c_0] = \Pr[C = c_1]$, where
  $C$ is the ciphertext corresponding to a message chosen according to
  $\mathcal{D}$.
  
  For this exercise, define the cipher space as the the set of
  ciphertexts that can be returned with nonzero probability for a
  uniformly random message (that is, the set of ciphertexts that can
  ever be returned by the encryption function).
\end{exercise}

\begin{answer}
  Solution here.
\end{answer}

% ====================================================================
% ====================================================================

\begin{exercise}
  \textbf{(20 points)} 
Consider the probability ensembles $\{X_n\}_n$ and 
$\{Y_n\}_n$. Consider distinguishers $\mathcal{D}_1, \ldots, \mathcal{D}_m$ where 
$m = \poly(n)$ 
for some polynomial function $\poly(\cdot)$, 
such that only one can distinguish $\{X_n\}_n$ and $\{Y_n\}_n$ 
with non-negligible probablity, but you do not know which one. 
Construct a distinguisher $\mathcal{D}^*$ that can distinguish between $\{X_n\}_n$ 
and $\{Y_n\}_n$ with non-negligible probability. 
Please describe your construction of $\mathcal{D}^*$ and provide a proof. 
\end{exercise}


\begin{answer}
  Solution here.

\end{answer}

% ====================================================================
% ====================================================================

\end{document}